\documentclass[10pt,a4paper]{amsart}
\usepackage[margin=19mm]{geometry}
\usepackage{amsmath,amssymb,mathtools}
\usepackage[scr=rsfs,cal=euler]{mathalfa}
\usepackage{xfrac}
\usepackage[unicode,hidelinks,bookmarks=false]{hyperref}
\newcommand{\ie}{i.e.\ }
\newcommand{\cf}{cf.\ }
\newcommand{\resp}{respectively\ }
\newcommand{\Subsection}{\subsection}
\providecommand{\nobreakdash}{\nobreak\hbox{-}}
\setlength{\emergencystretch}{2em}
\multlinegap0pt
\usepackage{bm}
\usepackage{bbm}
\usepackage{dsfont}
\usepackage{xspace}
\usepackage{cancel}
\usepackage{mathtools}
\usepackage{amsthm}
\makeatletter
\@ifundefined{lemma}{\newtheorem{lemma}{Lemma}}{}
\makeatother
\title[Local rigidity of convex hypersurfaces in spaces of constant]{Local rigidity of convex hypersurfaces in spaces of constant curvature}
\author{Alexander A. Borisenko}
\date{Source: arXiv:2505.18564v2 (CC BY 4.0)}
\begin{document}
\maketitle

\noindent\emph{Source and licence.} Local rigidity of convex hypersurfaces in spaces of constant curvature. Version arXiv:2505.18564v2 (CC BY 4.0), distributed under the Creative Commons Attribution 4.0 International licence, as recorded in the arXiv licence metadata for this exact version and shown on the abstract page https://arxiv.org/abs/2505.18564v2. Full rights notice: licenses/arxiv-2505.18564-CC-BY-4.0.txt.\par\smallskip

\noindent\textit{Editorial scope.} Counted unit: the Lemma whose source label is \texttt{Lem:3} in the article (source0.tex lines 161--164, source label \texttt{Lem:3}), delivered together with the article's own complete original proof of it (source0.tex lines 166--219, one \texttt{proof} environment of 54 lines). No mathematics is written, repaired or re-derived: statement and proof are the article's own text line for line. Required context retained uncounted: none. Every definition, display and label of those lines is retained unchanged. Omitted from this package and not redistributed: 1--160, 165--165, 220--417. Nothing inside a retained excerpt is omitted or altered: every retained body is delivered exactly as the article's lines are written, so no omission receipt of any kind accompanies this package. Citations: the retained excerpts carry no citation command at all, so no bibliography entry of the article is transcribed and no reference is redistributed. Reference adaptation: none was needed and none was made; every reference inside the delivered unit resolves inside it, so no unresolved reference or citation is produced. Printed slips kept literal: no printed word, symbol, hypothesis or number of the counted statement or of its proof was altered. Layout and class: the article's own class is \texttt{amsart} with packages including amsmath, geometry, hyperref, graphicx; the wrapper is the lane's pinned \texttt{amsart} editorial wrapper at 10pt on A4, which restates the theorem-like environments the retained text needs and adds the article's own macro definitions that the retained text needs (none), with extra wrapper packages (bm, bbm, dsfont, xspace, cancel, mathtools). The delivered numbering is the unit's own. Assets: no retained span contains a figure environment, a graphics-inclusion command, a PDF-page inclusion, a table or a TikZ picture; no image file is copied into this package, the package declares no asset dependency and no image file is redistributed with it. Licence: version arXiv:2505.18564v2 of this article is distributed under the Creative Commons Attribution 4.0 International licence, as recorded in the arXiv licence metadata for this exact version and shown on the abstract page https://arxiv.org/abs/2505.18564v2. A full rights notice is delivered at licenses/arxiv-2505.18564-CC-BY-4.0.txt. Characters: every retained excerpt is pure ASCII, so the delivered engine, XeLaTeX with its default Latin Modern text font, reports no missing character.
\begin{lemma}
\label{Lem:3}
Let $F_1, F_2$ be two closed nondegenerate convex curves with marked points $P_1, P_2$. Suppose that $F_1, F_2$ have the same orientation and the same length $s_0$. Then there exists an isometry of the plane that moves these curves so that the angles between the right semi-tangents at the points $P_1(s), P_2(s)$ is less than $\pi$, where $s$ is the arc length measured starting from $P_1=P_1(0), P_2=P_2(0)$. 
\end{lemma}

\begin{proof}
Let us align the points $P_1(0)$ and $P_2(0)$ using an appropriate isometry of the plane in such a way that these points coincide and that their right semi-tangents at these points also coincide. 
If the curve $F_1$ is strictly convex, then it is possible to choose as the point $P_1(0)$ any point of $F_1$. If on the curve $F_1$ there exists a line segment, then the beginning of the segment we choose as the point $P_1(0)$. The orientation on the segment is the same as the orientation on the curve $F_1$.
If for every $s \in (0, s_0]$, the angle between the right semi-tangent at $P_1(s)$ and $P_2(s)$ is $<\pi$, then we are done. So suppose this is not the case. Then there is a pair of points $P_1(\sigma_0)$, $P_2(\sigma_0)$ such that the right semi-tangents at those points have the opposite direction. For definiteness, we can assume that the turning of the arc $P_1(0)P_1(\sigma_0)$ is bigger than the turning of the arc $P_2(0)P_2(\sigma_0)$. Then
\[
\tau_1(\sigma_0) = \tau_2(\sigma_0) + \pi,
\] 
where for $i \in \{1,2\}$, $\tau_i(s)$ is the turning of the arc of the curve $F_i$ from the point $P_i(0)$ to the point $P_i(s)$. Observe that $\tau_2(\sigma_0) < \pi$ because the complete turning of a closed convex curve is equal to $2\pi$ and the point $P_1(\sigma_0)$ does not coincide with the point $P_1(0)$. Furthermore, both $\tau_1(s)$ and $\tau_2(s)$ are non-negative and non-decreasing functions.

For each $s \in [0, s_0]$, let us align the points $P_1(s)$ and $P_2(s)$ in the same way as we did above for $P_1(0)$ and $P_2(0)$, and assume that we have the same situation with the oppositely oriented right semi-tangents. That is, for the points $P_1(s)$, $P_2(s)$ we would find the points $P_1\big(f(s)\big)$, $P_2\big(f(s)\big)$ such that
\[
\omega_1(s) = \omega_2(s) + \pi, \quad 0 < \omega_2(s) < \pi,
\]  
where $\omega_1(s)$, $\omega_2(s)$ are the turnings of the arcs $P_1(s) P_1\big(f(s)\big)$, $P_2(s) P_2\big(f(s)\big)$.

Let us pick such $s$ so that
 \begin{equation}\label{turning}
    2\pi > \tau_2(s) + \omega_2\big(f(s)\big) = \tau_2\big(f(s)\big) > \pi.
 \end{equation}
 
Then $\tau_1\big(f(s)\big) > 2\pi$. But since $\tau_2\big(f(s)\big) < 2\pi$, then $f(s) < s_0$ and $\tau_1\big(f(s)\big) < 2\pi$. This is a contradiction. Therefore, there exists $\sigma_0 < s_0$ such that the isometry that maps $P_1(\sigma_0)$ to $P_2(\sigma_0)$, and the right semi-tangent at $P_1(\sigma_0)$ to the right semi-tangent at $P_2(\sigma_0)$ is the required isometry.

Let us explain it in detail. Let us take on $F_2$ the point $P_2(s)$ such that the turning 
$\tau(P_2(0)P_2(s))=\pi$. It is possible the possible cases: \\
1) such point $P_2(s)$ exists. If $\omega_2(s) >0$, then the inequality (\ref{turning}) is true; \\
2) Let us assume that $\omega_2(s)=0$. Then 
$$ 2\pi > \tau_2(s)+\omega_2\big(f(s)\big)=\pi \quad \text{and} \quad \tau_1\big(f(s)\big)=2\pi.$$
But $\tau_2\big(f(s)\big)= \pi$ and it follows that $f(s) < s_0$. The point $P_1(f(s))$ does not coincide
with  the point $P_1\big(f(s)\big)$, and the segment $P_1\big(f(s)\big)P_1(0)$ belongs to the curve $F_1$. 
If the point $P_1(0)$ is a corner point, then turning of the curve $F_1 >2\pi$. 
It is impossible for a closed convex curve. 
If the point $P_1(0)$ is a smooth point, then we have a contradiction with the choice of the point $P_1(0)$. \\
3) The turning of the arc $P_2(0)P_2(s)$ is greater than $\pi$ in the case when the point $P_2(s)$ is a corner point. The turning of the arc $P_2(s-\varepsilon)P_2(s)$ goes to $\pi - \alpha$ if $\varepsilon \to 0$, $\alpha$ is the angle at the point $P_2(s)$. 
Suppose that on the arc $P_2(s-\varepsilon)P_2(f(s-\varepsilon))$ , $f(s-\varepsilon) >s$ the following holds
$$\omega_1(s-\varepsilon) = \omega_2(s-\varepsilon) + \pi.$$
Then on the arc $P_1(s-\varepsilon)P_1(s)$ there exists a point with parallel support lines with the 
distance $\leq \varepsilon$. The curves $F_1$, $F_2$ are nongenerate and for $\varepsilon < c$, 
where $c >0$ is a constant, it is impossible- 
Then for $\varepsilon < c$ it is true strong inequalities
$$ 2\pi > \tau_2(s-\varepsilon) +  \omega_2(s-\varepsilon) =   \tau_2\big(f(s-\varepsilon)\big) > \pi,  $$
and we obtain contradiction as in the case 1).



%%Let $\ell_1(\phi)$ and $\ell_2(\phi+\pi)$ be the supporting parallel lines of the curve $F_1$, where $\phi$ is the angle between a directed line and the $x$-axis on the plane with fixed orientation. We can find an angle $\phi_0$ such that the points of support of the supporting lines $\ell_1(\phi_0)$ and $\ell_2(\phi_0 + \pi)$ divide the curve $F_1$ into two arcs of equal length. Denote $P_1(0)$ and $P_1(s_0/2)$ those points of support. Move the curve $F_2$ such that the corresponding under the isometry point $P_2(0) \in F_2$ coincides with $P_1(0)$ and the right semi-tangents at those points also coincide.
%%
%%Suppose that $P_1(s_0/2 + \sigma_0)$, $P_2(s_0/2+\sigma_0)$ be the first points along the curves such that the angle between the corresponding right semi-tangents is equal to $\pi$. Then the turning $\tau$ of the arc $P_2(0)P_2(s_0/2 + \sigma_0)$ will be less then $\pi$, while the turning of the arc $P_1(0) P_1(s_0/2 + \sigma_0)$ is equal to $\pi+\tau$.
%%
%%Let $P_2(s_0/2 + \sigma_1)$, $\sigma_1 > \sigma_0$, be a point such that at this point the supporting line is parallel to the right semi-tangent of the curve $F_2$ at the point $P_2(0)$. Let us rotate the arc of $F_1$ clockwise in such a way so that the point $P_1(s_0/2 + \sigma_1)$ aligns with the point $P_2(s_0/2 + \sigma_1)$ and their supporting lines align too. Then the points $P_1(s_0 + \sigma_1)$, $P_2(s_0 + \sigma_1)$ be the new points with respect to which the arc length parameter will be measured $P_1'(0) = P_2'(0)$. Observe that the orientation of curves $F_1, F_2$ was chosen to be counter clockwise.
%%
%%The turning of the arc $P_1(s_0/2 + \sigma_1)P_1(s_0/2)$ is equal to $\tau_1 < \tau$, the turning of the arc $P_1(s_0/2+\sigma_0)P_1(s_0/2)$ is equal to $\tau_0 < \tau_1$. Let us measure the corresponding under the isometry points starting with the point $P'_1(s_0/2+\sigma_1) = P_2(s_0/2+\sigma_1)$. We will move in the clockwise direction. We denote with prime the corresponding points on $F_1$ after the rotation. The turning of the arc $P_1'(s_0/2+\sigma)P_1'(s_0/2)$ is equal to $\tau_1 < \pi$. The turning of the arc $P_1'(s_0/2+\sigma_1)P_1'(0)$ is equal to $\pi + \tau_1$. Therefore, in the arc $P_1'(s_0/2)P_1'(0)$ there are no points with the opposite direction of the supporting semi-tangents. This is because the turning of the curve $P_1'(s_0/2)P_1'(0)$ is equal to $\pi$, while the turning of the corresponding arc on the curve $F_2$ is equal to $\tau_0 < \tau_1 < \pi$. The turning of the curve $P_2(0)P_2(s_0/2 + \sigma_1)$ is equal to $\pi$, and the turning of the corresponding isometric arc is $\pi - \tau_1$.
%%
%%Since the arcs $P_2(0)P_2(s_0/2 + \sigma_1)$, $P_1'(0)P_1'(s_0/2+\sigma_1)$ are located in such a way that the difference between the turnings of the arcs $P_2(0)P_2(\sigma)$, $P_1'(0)P_1'(\sigma)$ is always less than $\pi$, we conclude that these arcs there will be no points at which the corresponding under the isometry semi-tangents have the opposite direction. 
\end{proof}

\end{document}
